\section{Supplementary Evaluation Details}~\label{appx:exp}

This appendix specifies the threshold protocol and the failure categories used in Sec.~\ref{exp:sys_baseline}. The reported outcomes use the same trial sets as the corresponding main-text comparisons.

\subsection{Threshold Protocol}~\label{appx:threshold_exp}
The threshold study evaluates $\tau\in\{0.0,0.1,\ldots,1.0\}$ with entropy gain query selection and oracle responses. Each threshold has $200$ trials per domain. The planner parameters are those in Appendix~\ref{appx:planner_params}, and the shared reference at $\tau=0.8$ is specified in Appendix~\ref{appx:oracle_reference}. A threshold of zero suppresses queries because belief confidence cannot fall below zero. A threshold of one requires confidence to reach one before the robot continues, unless no differentiating fact remains. Table~\ref{tab:threshold_results} reports the success rates and successful-trial query counts.

We use $\tau=0.8$ as a common operating point for the two domains. In the reported trials it gives $100\%$ success in WasteSorting and $98\%$ in TomatoHarvesting while requiring fewer queries than $\tau=1.0$. This empirical success--query trade-off does not establish a universally optimal threshold or an independently validated choice.
% TODO(author): Confirm whether tau=0.8 was fixed before the sweep or selected
% after inspecting it. The repository report describes a trade-off choice,
% but does not document a dated selection decision or a held-out protocol.

\subsection{Ablation Reporting}~\label{appx:policy_comparison}
Table~\ref{tab:policy_ablation} contains Fully-Random, W1--W3, Q1--Q2, and Ours. Each configuration has $200$ attempted trials per domain. The success denominator includes the four unparseable W2 outputs. Cost metrics include successful trials only. The policy definitions and domain-specific discount factors are given in Appendix~\ref{appx:ablation_configs}.

\subsection{VLM Failure Classification}~\label{appx:vlm_failure_analysis}
We classify failed physical trials using the available action candidates, external responses, and subsequent execution. A VLM failure occurs if a response incorrectly verifies a state fact or selects an action that causes failure despite an available feasible alternative. Feasible actions that do not advance the task can also cause a VLM failure if the resulting sequence exhausts the trial limit. A planner failure occurs if the candidate set contains only infeasible actions or the baseline autonomously selects an action whose preconditions are not satisfied. The termination label alone does not determine the category.

Table~\ref{tab:app_vlm_failures} reports the failure counts. In WasteSorting, Ours has three VLM feedback failures. KnowNo has three VLM failures and $17$ planner failures. In TomatoHarvesting, Ours has three VLM feedback failures. KnowNo has $18$ VLM failures and $12$ planner failures. Six of the $18$ VLM failures involve detection choices that exhaust the step budget despite feasible alternatives for task progress. The remaining $12$ involve incorrect action selections.

\begin{table}[!htbp]
\centering
\caption{Failure counts in the VLM evaluation. Percentages use failed trials as the denominator.}
\label{tab:app_vlm_failures}
\small
\begin{tabular}{llrrr}
\toprule
Domain & Method & Failed trials & VLM failures & Planner failures \\
\midrule
WasteSorting & KnowNo & $20$ & $3$ ($15\%$) & $17$ ($85\%$) \\
WasteSorting & Ours & $3$ & $3$ ($100\%$) & $0$ \\
TomatoHarvesting & KnowNo & $30$ & $18$ ($60\%$) & $12$ ($40\%$) \\
TomatoHarvesting & Ours & $3$ & $3$ ($100\%$) & $0$ \\
\bottomrule
\end{tabular}
\end{table}

\subsection{KnowNo Action Requests}~\label{appx:knowno_impl}
KnowNo generates candidate actions from domain information and the current symbolic state. Candidate scores and the calibrated threshold define the conformal prediction set. If the set contains multiple actions, the interface requests one action choice from the external expert. KnowNo requests one response per trigger. The selected action is then executed, and the execution result updates the state information used for the next decision. Appendix~\ref{appx:baseline_configs} lists the model and calibration parameters.
